\chapter{Theory}

This chapter builds the theoretical case for the approach taken in this thesis.
The argument unfolds in three movements: first, why large language models alone are
insufficient for multi-document legal case analysis; second, why retrieval-augmented
generation --- the standard remedy --- carries fundamental limitations that make it
inadequate for cross-document reasoning; and third, what the literature suggests as
more capable alternatives, which together motivate the structured extraction approach
evaluated in this thesis.

% -------------------------------------------------------
\section{Large Language Models and Context Windows}
% -------------------------------------------------------

Modern LLMs are built on the transformer architecture \cite{vaswani2023attentionneed},
which processes input as a sequence of tokens bounded by a fixed context window.
Recent models advertise increasingly large context windows --- some exceeding one
million tokens --- which may appear to solve the problem of large document corpora.
This appearance is misleading for at least three reasons.

First, the nominal context size overstates effective capacity. Hsieh et al.\ introduce
RULER, a benchmark measuring the \textit{real} usable context of long-context LLMs,
and demonstrate that performance degrades substantially before the advertised limit is
reached \cite{hsieh2024ruler}. Second, LLMs do not attend uniformly across the context:
performance is highest when relevant information is at the beginning or end, and
degrades for information buried in the middle --- the \textit{lost in the middle}
phenomenon \cite{liu2023lostmiddlelanguagemodels}. Third, and most practically, a
legal case involving dozens of attachments spanning hundreds of pages cannot be loaded
in full for every query without prohibitive cost and latency, regardless of model
capability.

The implication is a design principle: rather than maximizing the amount of text in
context, the goal should be to include the \emph{minimum} amount of information while
ensuring that information is as relevant as possible to the query at hand. This principle
rules out naive full-document injection and motivates retrieval as the primary mechanism
for selecting what enters context. The question then becomes: what are the limits of
retrieval?

% -------------------------------------------------------
\section{Retrieval-Augmented Generation}
% -------------------------------------------------------

Retrieval-Augmented Generation (RAG) operationalizes the minimum-relevant-information
principle by retrieving only the most semantically similar document chunks at query time,
rather than loading entire documents \cite{lewis2021retrievalaugmentedgenerationknowledgeintensivenlp}.
Documents are chunked and embedded offline into a vector store; at inference time, the
top-$k$ most similar chunks are retrieved and inserted into the LLM context alongside
the query.

For straightforward factual lookup --- ``what is the defendant's registered address?''
--- this works well. The answer lives in a single location in the document corpus, it
is semantically similar to the query, and a chunk-based retriever will find it. The
problems emerge when queries require understanding that spans multiple documents, or
when the relevant information is not directly expressible as a similarity to the query
string itself.

% -------------------------------------------------------
\section{Limitations of Naive RAG for Complex Reasoning}
\label{sec:rag-limitations}
% -------------------------------------------------------

\subsection{Semantic Similarity Does Not Guarantee Relevance}

The retrieval step selects chunks based on embedding similarity to the query. High
similarity does not guarantee true relevance: documents may share keywords or surface
features without addressing the underlying intent or logic of the query
\cite{gao2025synergizingragreasoningsystematic}. In legal contexts, where questions
about causation, liability, or damages require understanding relational chains across
documents, surface-level matching is an insufficient retrieval criterion.

Compounding this, injecting all retrieved chunks directly into context can produce
fragmented or contradictory input that confuses the LLM and degrades answer quality
\cite{gao2025synergizingragreasoningsystematic}. The problem is not just what is
retrieved but how it combines in context.

\subsection{Cross-Document and Multi-Hop Reasoning}

The more fundamental limitation is that the answer to a legal query often does not
reside in one or two isolated chunks, but is distributed across multiple documents.
Existing naive RAG systems are inadequate for queries that require retrieval and
reasoning over multiple pieces of supporting evidence across documents
\cite{tang2024multihoprag}. In multi-document settings, relevant information is spread
across sources and requires coherent synthesis: naively concatenating retrieved passages
leads to fragmented or contradictory responses, a problem that is especially acute in
legal and domain-specific QA where multi-hop reasoning is critical
\cite{liang2025reasoningrag12}.

This limitation is structural, not a matter of tuning. A chunk embedding represents a
local fragment of a single document. No similarity threshold or retrieval hyperparameter
can make a chunk-based retriever produce a coherent synthesis of facts distributed across
an entire case file.

\subsection{Unstructured Retrieval and Missing Information}

Standard RAG processes unstructured text passages without leveraging structured domain
knowledge, and struggles with complex queries requiring multi-hop reasoning or
domain-specific knowledge organization \cite{Jiang_2025}. Unstructured retrieval also
risks returning passages that contain non-atomic information --- facts entangled with
irrelevant surrounding text --- that can mislead the LLM rather than inform it
\cite{Jiang_2025}. Chunk size amplifies this: small chunks provide insufficient
context for understanding a passage, while large chunks dilute the semantic signal and
retrieve off-topic content.

These three limitations --- relevance mismatch, cross-document reasoning failure, and
unstructured knowledge gaps --- together define the core challenge this thesis addresses.
A legal agent that relies solely on chunk-based RAG will, by design, struggle with
case-level questions that require a coherent picture of the whole. This motivates
approaches that move beyond isolated chunk retrieval toward structured, relational
representations of case knowledge.

% -------------------------------------------------------
\section{Graph-Based and Multi-Hop Retrieval}
% -------------------------------------------------------

One line of work addresses RAG's limitations by enriching retrieval with relational
structure. GraphRAG \cite{edge2025localglobalgraphrag} builds a knowledge graph over
the document corpus, enabling community-level summarization and structured traversal of
entity relationships rather than relying on chunk similarity alone. HopRAG
\cite{liu2025hoprag} extends this with explicit multi-hop reasoning, constructing
logical chains across documents to answer queries that require synthesizing evidence
from multiple sources.

These approaches directly target the cross-document reasoning failure described above.
Rather than asking whether a chunk is similar to a query, they ask how entities and
facts across the corpus are related. The factsheet approach implemented in this thesis
can be understood as a lightweight, domain-specific variant of this idea: rather than
constructing a general-purpose knowledge graph, it extracts and organizes precisely the
information most relevant to legal case management --- parties, claims, events, damages,
and governing law --- into a structured, queryable representation that supplements RAG
for case-level queries.

% -------------------------------------------------------
\section{Agentic Context Management}
% -------------------------------------------------------

The final theoretical component concerns how structured knowledge is maintained over
the lifecycle of a multi-turn conversation. Rolling summarization --- periodically
condensing conversation history into a compact summary --- is a practical strategy for
preventing context overflow in long sessions \cite{microsoft2024semantickernelpython}.
Stateful agent frameworks such as LangGraph enable persistent state management through
checkpointing, allowing an agent to maintain structured representations that accumulate
case knowledge across turns without requiring the full document history to remain in
context \cite{langgraph}.

In the system developed in this thesis, the factsheet is the primary persistent state
artifact. It is initialized during document ingestion and updated incrementally as new
attachments are processed and new facts confirmed, providing the agent with a structured,
token-efficient overview of the case at every point in the conversation. The theoretical
argument, then, is that this persistent structure addresses what RAG cannot: a coherent,
relational picture of the case that is available without vector-based retrieval: the factsheet is loaded directly from the structured database, not assembled from semantically matched chunks at query time.

The following chapter describes how these theoretical motivations are translated into
system design and how their combined effect is evaluated empirically.

% -------------------------------------------------------
\section{Related Work}
% -------------------------------------------------------

This section surveys prior work on the core problems this thesis addresses: overcoming RAG limitations in large multi-document corpora, structured information extraction for complex QA, and agentic LLM systems for professional knowledge work. The legal domain is the application context, but the architectural challenges --- handling large, heterogeneous document collections with distributed factual content --- are not unique to law. Related problems appear in financial due diligence, medical records analysis, and other high-stakes professional domains where information is voluminous, structured, and spread across many sources.

It is worth noting that most prior work on LLMs in legal settings concerns legal \emph{reasoning} --- statutory interpretation, case law analysis, and normative prediction \cite{chalkidis2020legalbert, guha2023legalbench}. This thesis is not primarily concerned with legal reasoning ability. The focus is on the prior, more practical problem: how to make the factual content of a large case file reliably accessible to an LLM agent, regardless of which model underlies it. This distinction matters when interpreting the related work: findings about legal language models and legal reasoning benchmarks address a different problem, and this thesis fills a gap that those works do not cover.

\subsection{Retrieval-Augmented Generation for Legal QA}

RAG has emerged as the standard architecture for domain-specific document QA, and its application to legal settings has been explored in several directions. Louis et al.\ \cite{louis2024interpretable} demonstrate that RAG-augmented LLMs achieve competitive performance on long-form legal question answering while producing interpretable, citation-grounded responses. Their system retrieves relevant statutory passages and case law and incorporates them as context at inference time --- the same basic mechanism used by the baseline agent in this thesis. Their work highlights both the strengths of RAG (traceable sourcing, factual grounding) and its limitations when queries require synthesising evidence from across a large corpus.

The fundamental tension identified in RAG-based legal systems is that legal questions rarely decompose into a single, retrievable passage. Multi-hop reasoning over distributed evidence is common: establishing liability may require linking a party's obligation (in a contract), its breach (in correspondence), and the resulting damage (in a financial report), none of which are co-located in the document corpus. MultiHopRAG \cite{tang2024multihoprag} and HopRAG \cite{liu2025hoprag} address this limitation in general QA settings by explicitly chaining retrieval steps, constructing logical paths across sources. This thesis takes a complementary approach: rather than chaining retrievals reactively at query time, it extracts and pre-structures the key relational facts during document ingestion, making them available as a persistent case state without requiring multi-hop retrieval.

\subsection{Structured Extraction and Knowledge Organisation}

Extracting structured information from unstructured legal documents --- parties, obligations, events, claims, and temporal relationships --- has been studied as an information extraction problem long before the LLM era. Recent work has revisited this problem with LLMs: Jiang et al.\ \cite{Jiang_2025} propose Retrieval and Structuring Augmented Generation (RSAG), which combines retrieval with LLM-driven structuring of retrieved content before generation. Their results show that structuring retrieved passages into typed entities before injecting them into the generation context consistently outperforms naive RAG on complex reasoning tasks. This provides independent empirical support for the core design hypothesis of this thesis: that structured, relational case state is more useful to an LLM than raw retrieved chunks.

GraphRAG \cite{edge2025localglobalgraphrag} extends this idea to the full document corpus, building a knowledge graph over entities and their relationships to enable community-level summarization and structured traversal. The factsheet implemented in this thesis can be understood as a lightweight, domain-specific instance of the same principle: rather than constructing a general-purpose property graph, it organises only the entities most relevant to legal case management into a typed, relational structure. The practical trade-off is that a domain-specific schema can be designed to match exactly the queries a legal practitioner asks, at the cost of being less general than a full knowledge graph.

\subsection{Agentic Memory Management}

A structural constraint shared by all retrieval-based approaches is that the LLM context window is finite. MemGPT \cite{packer2024memgptllmsoperatingsystems} addresses this directly by treating the context window as a fixed-capacity ``main context'' --- analogous to RAM --- and introducing a tiered memory architecture in which an explicit memory manager decides what information to page in from external storage and what to evict. Rather than relying on a single retrieval query, the system supports fine-grained memory operations: appending new facts to context, archiving stale content to external storage, and selectively restoring it on demand. The result is an agent that can maintain coherent state across long interactions without ever exceeding the physical context limit.

Several of these mechanisms inform the design of the system in this thesis. The factsheet functions as a domain-specific persistent memory: it is initialised during document ingestion and injected directly into context at query time, bypassing the need for vector retrieval of case-level facts. Rolling summarisation of conversation history serves the eviction function, preventing the accumulated dialogue from consuming the context budget. What the factsheet approach does not yet implement --- and what MemGPT handles explicitly --- is active monitoring of context utilisation with tool-mediated relief: notifying the model when capacity is near and providing tools to archive or compact in-context content. This remains a natural direction for future work.

A complementary diagnostic perspective is offered by Yuan et al.\ \cite{yuan2026diagnosingretrievalvsutilization}, who decompose agent memory failures into two distinct bottlenecks: \emph{retrieval failures}, in which the correct information is not surfaced from storage, and \emph{utilisation failures}, in which the correct information is retrieved but not correctly used by the model. Their empirical analysis shows that improving retrieval quality consistently yields larger performance gains than increasing the sophistication of the write-time encoding process. This finding supports the design priority in this thesis: the factsheet eliminates retrieval entirely for case-level facts by placing them directly in context, targeting the retrieval bottleneck rather than investing in more elaborate embedding strategies.

\subsection{Hierarchical and Abstractive Document Summarisation}

RAPTOR \cite{sarthi2024raptorrecursiveabstractiveprocessing} takes a different approach to the multi-document comprehension problem. Rather than retrieving isolated chunks, it recursively clusters document segments and generates abstractive summaries at each level of the cluster hierarchy, constructing a tree whose leaves are raw text chunks and whose internal nodes are progressively higher-level summaries. At query time, retrieval can target any level of the tree, enabling both fine-grained lookup and broad, thematic search that spans many documents.

RAPTOR directly addresses the cross-document reasoning failure identified in Section~\ref{sec:rag-limitations}: because higher-level nodes aggregate content from many source segments, a single retrieval step can surface information that would otherwise require multiple hops across the corpus. In practice, however, the approach carries significant costs. Constructing the hierarchy requires summarising every cluster at every level, which is computationally expensive and heavily dependent on the quality of the summarisation model. Weak summarisation propagates errors upward through the tree, and the cost scales with both corpus size and model capability. Furthermore, the summaries produced by RAPTOR are consumed solely as retrieval targets: they improve the chances that a relevant passage is found, but they are not themselves injected as structured context.

The factsheet approach taken in this thesis can be understood as a more targeted variant of the same intuition. Rather than abstracting the entire corpus into a general-purpose hierarchy, it extracts only the entities most relevant to legal case management --- parties, claims, damages, obligations, and key events --- and organises them into a typed, relational structure. This structure serves a dual purpose: it both enables structured retrieval of specific case facts and is injected directly as context, giving the model a coherent case overview without any retrieval step. The trade-off is scope: the factsheet covers only what a legal practitioner is likely to ask, rather than supporting arbitrary queries over the full document space.

\subsection{Research Gap}
\label{sec:research-gap}

Taken together, the related work establishes several important results: (1) general RAG is insufficient for multi-document legal reasoning due to relevance mismatch and cross-document synthesis failures; (2) structured extraction consistently improves performance on complex QA tasks; (3) graph-based and multi-hop retrieval partially address cross-document reasoning by making entity relationships explicit; (4) tiered memory management, as demonstrated by MemGPT, enables coherent long-horizon agent behaviour by selectively controlling what enters context; and (5) hierarchical abstractive summarisation, as in RAPTOR, improves broad-coverage retrieval but at significant computational cost and without producing a directly injectable, domain-typed representation.

What none of these lines of work address is the combination: a domain-specific, persistently maintained structured case state that is injected directly into context --- bypassing retrieval for the information that matters most --- within a conversational agent evaluated on authentic legal case files across multi-session workflows. Prior systems are either domain-agnostic (MemGPT, RAPTOR, GraphRAG), evaluated on synthetic or curated benchmarks rather than real case management tasks, or focused on legal reasoning rather than the practical problem of making large, heterogeneous case files reliably accessible. This thesis addresses that gap by implementing and empirically evaluating this architecture against a RAG baseline using real legal case documents and expert-constructed ground truth.

